% !TEX root = ../main.tex
\chapter{Introduction and Research Problem}
\label{ch:introduction}

Decentralized Finance (DeFi) transformed how capital is allocated on public blockchains: lending, borrowing, and leveraged trading can operate without traditional intermediaries \cend{1}. However, most credit-like transactions remain highly collateralized since protocols cannot depend on off-chain identity frameworks or existing credit data \cend{2}. Within the context of pseudonymity, the key issue is how to efficiently distribute capital among unknown parties.

This work tackles a specific research question: can an on-chain reputation metric derived from ERC-20 \texttt{approve}/\texttt{permit} events, called EndorseRank, be calculated less expensively than transfer-based Adaptive Weighted PageRank (AWP) on Arbitrum One, does it capture endorsements more accurately, and does it deviate from AWP in a verifiable manner? The core evaluation compares EndorseRank to AWP on the same set of wallets, validated using matching transfer and approval statistics and a holdout period of future approvals, following the rank-correlation testing method described with AWP\cend{3}. In this paper, an on-chain reputation metric is defined as a value derived from on-chain activity and interactions without any off-chain identity information \cend{4}, and EndorseRank is that value calculated from approval edges.

\dissertationSubheading[sec:background]{Background: DeFi, collateralization, and pseudonymity}

Public blockchains enable programmable settlement without a central authority \cend{5}. Terms of a deal can be encoded into a smart contract \cend{6}, and a DeFi platform that uses such a contract clears trades on the blockchain instead of through a custodian \cend{7}. DeFi protocols subsequently stack these contracts together to build decentralized markets for lending, automated market making and perpetual futures \cend{8}. On Arbitrum One and other L2 networks, there is sufficient activity to inspect transfers, approvals, and protocol events at the level of individual wallets.

However, despite this visibility, DeFi credit markets remain inherently over-collateralized. Most protocols require borrowers or traders to supply collateral greater than the size of their position, since trading partners are pseudonymous accounts rather than verified individuals \cend{9,10}. Pseudonymity is an inherent property of public ledgers: anyone can generate addresses for almost no cost, and one entity can own multiple wallets \cend{11}. Conventional credit rating relies on personal histories that do not easily map onto this environment \cend{2}. Therefore, risk control relies on collateral reserves, liquidation mechanisms, and protocol rules, rather than personalized reputation.

Collateralization safeguards protocols, but requires inefficient use of capital. Traders and borrowers are forced to tie up funds they might otherwise invest; platforms do not offer under-collateralized goods that traditional finance provides to those with a credit record. Researchers have therefore looked at on-chain signals as a proxy for off-chain identity: transfer-graph centrality \cend{4}, network risk propagation \cend{12}, prospect theory and blockchain-based credit scoring \cend{13}, and machine learning pipelines built over lending protocols \cend{14}. All of these share a common assumption: that publicly available ledger data contains information that can be used to assess risk while avoiding the issue of deanonymizing users.

Graph ranking algorithms appear to be well-suited to this problem. PageRank and its variants compute rankings of nodes using recursive endorsement, where each node's score is computed based on endorsements from other nodes \cend{15}. Prior web trust algorithms such as EigenTrust \cend{16} and TrustRank \cend{17} use similar approaches to assess peer-to-peer and web-spam trustworthiness. Do and Do\cend{4} use both PageRank and HodgeRank to assess the trustworthiness of Ethereum users using a graph constructed from transfers, and call this a form of social credit. Adaptive Weighted PageRank (AWP), proposed by Do, Do, and Nguyen\cend{3}, is another transfer graph ranking algorithm that weights edges based on transfer values and applies logistic time decay, and they demonstrate the validity of their algorithm by showing that it correlates with inbound transfer count and value as proxies for economic importance. This body of work shows that transfer graph reputation forms a credible baseline. However, it does not take into account explicit spending authorization, a separate on-chain action recorded as ERC-20 allowances \cend{18,19}.

\dissertationSubheading[sec:motivation]{Motivation for on-chain reputation}

There are four main factors that drive the need for on-chain reputation research in the context of leveraged DeFi.
\begin{itemize}
\item Capital efficiency: protocols are looking for ordinal ranking signals that could be used for tiered services (for example, more frequent limit increases, lower fees, or higher priority in monitoring), and perhaps even for differentiated collateral requirements in future controlled experiments, while maintaining on-chain settlement \cend{1}.
\item Auditability: unlike off-chain credit scores which are opaque, graph-based scores computed from public event data can be recomputed and audited, which would help alleviate some of the concerns regarding black-box decentralized credit scoring \cend{2}.
\item Operational cost: if the reputation system needs to be updated often (for example, before each leverage adjustment), it cannot afford to aggregate historical transfers every time it runs.
\item Standardization: there is no standard ordinal language for counterparty differentiation at the wallet level in pseudonymous DeFi, whereas screening ratios serve this purpose in other markets. Chapter~\ref{ch:discussion} explores potential use cases for industry adoption, where EndorseRank and AWP can provide infrastructure that allows protocols to offer differentiated fees, utilization, and monitoring efficiency, without replacing collateral.
\end{itemize}

Transfer-based AWP meets the first two criteria but is expensive for the third. To construct $G_{\text{transfer}}$, one must aggregate the user's transfer history, apply logistic time decay to these edges, and then run PageRank on this dense set of edges \cend{3}. Allowance graphs represent a different structural alternative. Whenever a user approves an ERC-20 token, the contract records how much of that token a spender is authorized to transfer \cend{18,19}. By looking at the most recent allowance for each owner--spender pair, one gets a snapshot of current spending authorization: since any new approval will supersede older approvals, there is no need for temporal decay. Since approvals are less frequent than transfers, the endorsement graph $G_{\text{approve}}$ on the same wallets has an order of magnitude fewer edges than the transfer graph $G_{\text{transfer}}$, and PageRank, which costs $O(|E|)$ per iteration, finishes proportionally faster (Table~\ref{tab:benchmark-runtime}). The speedup is an artifact of the sparser graph under latest-allowance snapshots and not an algorithmic improvement, as the same solver is used in both cases. Claim~C1 includes it as an operational property of the representation.

Beyond its efficiency, allowances have different semantics. A transfer could be for routing, market making, or a passive receipt. In contrast, an allowance is an active grant of permission to spend funds. By representing allowances as endorsement edges, the citation analogy of PageRank \cend{15} is mapped onto a trust relationship that protocols already depend upon when considering token interactions. It remains to be seen whether this semantic distinction yields rankings that are different, and if so, whether the rankings correlate differently with authorization and transfer checks. We address these questions empirically on Arbitrum One with a matched sample of wallets that have non-zero endorsement- and transfer-graph connectivity ($n=5{,}521$).\footnote{The matched sample is wallets with non-zero connectivity in the endorsement and transfer subgraphs during the observation period. To assess runtime sensitivity to node count, we use an expanded sample of $N=99{,}080$ wallets.}

\dissertationSubheading[sec:research-problem]{Research problem}

Adaptive Weighted PageRank (AWP) methods use the history of transfers to infer trust using weighted sums and logistic time-decay \cend{3}. Though they provide a rich economic interpretation, AWP (i) requires aggregation of costly historical transfers, (ii) models a valuable wallet as one that receives a large number of inbound transfers, and (iii) ignores the role of spending authorization. We propose EndorseRank, which uses the latest allowance of an owner--spender pair as the edge strength of the directed endorsement graph. The graph is much sparser, the signal is different, and there is no need for temporal decay.

Our research question is as follows: on an identical sample of Arbitrum One wallets, do the wallet rankings produced by EndorseRank differ from those produced by AWP in terms of their construct validity, correlation with allowance and transfer checks, prediction of future endorsements during a time holdout, and computational cost? The main comparison is between EndorseRank and AWP. Building on our first research question, we ask: since both methods perform random walks on two edge types across the same set of addresses, can we mix the allowance and transfer edges on a node-wise basis with a fixed weight $\lambda$ to define a Coupled PageRank (C-PR), and does any interior point in this family rank wallets better than either of the end members?

Table~\ref{tab:primary-methods-preview} summarizes the primary comparison. Table~\ref{tab:hrq-claim-map} summarizes hypotheses (H1--H3), research questions (RQ1--RQ5), and empirical claims (C1--C5) from Chapter~\ref{ch:results}.

\begin{table}[htbp]
\centering
\caption{Primary comparison: EndorseRank versus AWP (matched wallet cohort, $n=5{,}521$).}
\label{tab:primary-methods-preview}
\fitwidth{%
\begin{tabular}{p{0.18\linewidth}p{0.28\linewidth}p{0.48\linewidth}}
\toprule
Method & Input graph & Role \\
\midrule
EndorseRank & Latest ERC-20 allowances ($G_{\text{approve}}$) & Primary construct: allowance delegation; target of this research \\
AWP & Time-decayed transfers ($G_{\text{transfer}}$) & Primary baseline: IEEE RIVF 2023 transfer-graph method \cend{3}; Do and Do\cend{4} is the earlier PageRank/HodgeRank study on Ethereum transfers \\
C-PR & Both layers, mixed per node with weight $\lambda$ & Coupled operator whose end points $\lambda=1$ and $\lambda=0$ are EndorseRank and AWP; tests whether the two edge types combine (RQ5). S-PR, a transfer walk restarted from EndorseRank, is its ablation \\
\bottomrule
\end{tabular}
}
\end{table}

We define key terms here to avoid any potential confusion:

\begin{itemize}
\item \emph{DeFi (Decentralized Finance):} Financial services such as lending and borrowing implemented using smart contracts on public blockchains, generally without a central operator managing custody or settlement \cend{1}. DeFi protocols are composed with open interfaces, so that token movements can be observed on chain.
\item \emph{Smart contract:} Self-executing program code running on a blockchain that executes an agreement without a trusted intermediary \cend{6}. In this paper, smart contracts are the source of the allowance and transfer logs.
\item \emph{Wallet / Address:} A blockchain account that stores funds and interacts with smart contracts; this is the object whose reputation we wish to rank \cend{7}. Wallets are identified by their addresses; there is no assumption or recovery of any off-chain identity.
\item \emph{On-chain reputation score:} A numerical value computed from on-chain data representing a wallet's relative position for financial activities under pseudonymity \cend{4}. Scores are ordinal, that is, for ranking purposes; the absolute scale has no meaning in terms of credit limits.
\item \emph{PageRank:} A graph ranking algorithm that distributes influence through a directed graph using a damping factor \cend{15}. EndorseRank is PageRank applied to the latest-allowance endorsement graph: edge weights are equal to the current allowance value for each owner--spender pair, and no temporal decay is applied. Adaptive Weighted PageRank (AWP) is PageRank applied to transfer graphs, where edge weights are the transfer value times a logistic time-decay \cend{3}.
\item \emph{Coupled PageRank (C-PR):} A single random walk that traverses both allowance and transfer layers. The transition out of a wallet mixes the two row vectors using weights $\lambda$ and $1-\lambda$, then renormalizes over the layers the wallet actually occupies. $\lambda=1$ recovers EndorseRank and $\lambda=0$ recovers AWP on a common node set. Endorsement-seeded PageRank (S-PR) is the ablation that retains the transfer walk but restarts the process from the EndorseRank distribution, in the same way that TrustRank \cend{17} restarts from a set of seeds.
\item \emph{ERC-20 allowance / Approve / Permit:} Methods that keep track of how many tokens the owner authorizes the spender to spend on the owner's behalf \cend{18,19}. \texttt{Approve} is a regular on-chain call, and EIP-2612 \texttt{permit} allows the same state change through a signature. EndorseRank uses the latest non-zero allowance snapshot for each owner--spender pair.
\item \emph{Endorsement graph:} A directed graph in which edges denote an explicit delegation of spending authority rather than an actual transfer of assets \cend{18}. In EndorseRank, an edge $u \to v$ implies that wallet $u$ has approved wallet or contract $v$ to spend tokens belonging to $u$; the transfer-graph AWP is the opposite baseline \cend{3}.
\item \emph{Domain alignment:} The extent to which the reputation ranking is consistent with an observable proxy ranking, measured by Spearman's $\rho$ \cend{20} and Kendall's $\tau$ \cend{21}, in the construct-validity sense of Cronbach and Meehl\cend{22}. The alignment is construct specific: a method could be highly aligned with its own construct graph and lowly aligned with a different, orthogonal check.
\end{itemize}

\dissertationSubheading[sec:supplementary-extension]{Temporal holdout}

A contemporaneous allowance Kendall $\tau$ (Table~\ref{tab:alignment-allowance}) is an intended-construct check: PageRank on $G_{\text{approve}}$ should agree with the spender-side approval degree since the score is a global smoothing of this degree. To test whether the score conveys information beyond the same-window graph, the scores are fixed at 28~February~2026 and compared to two labels that only become available from March--May~2026: new owner--spender approval pairs (the endorsement construct), and new transfer senders (the flow construct). The holdout cohort is $t_1$ inbound spenders ($n=1{,}335$). Each PageRank is also compared against the raw degree it smooths at $t_1$, so that the gain of the propagated score over the local count is observed directly.

\begin{table}[htbp]
\centering
\caption{Summary of hypotheses, research questions, and empirical claims.}
\label{tab:hrq-claim-map}
\fitwidth{%
\begin{tabular}{p{0.12\linewidth}p{0.42\linewidth}p{0.40\linewidth}}
\toprule
Item & Addresses & Reported in \\
\midrule
H1 / RQ1 & Do EndorseRank and AWP yield distinct wallet orderings? & \S\ref{sec:rank-divergence}; inter-method $\tau$ with 95\% CI \\
H2 / RQ2 & Which social method aligns with allowance vs.\ transfer checks? & Tables~\ref{tab:alignment-family-ci} and~\ref{tab:tau-diff}; Claims C2--C3 \\
H3 / RQ3 & Is EndorseRank faster than AWP at equal $n$? & Tables~\ref{tab:benchmark-runtime}--\ref{tab:benchmark-scaling}; Claims C1, C4 \\
RQ4 & Do $t_1$ scores predict $t_2$ new approvals beyond same-window allowance degree? & \S\ref{sec:holdout-results}; Tables~\ref{tab:holdout-spenders} and~\ref{tab:holdout-diff} \\
RQ5 & Does the coupled operator C-PR beat both single-layer scores? & \S\ref{sec:coupled-results} and~\S\ref{sec:holdout-results}; Tables~\ref{tab:alignment-hybrid}, \ref{tab:tau-diff-hybrid} and~\ref{tab:holdout-diff} \\
C1 & EndorseRank vs.\ AWP runtime and graph size & Table~\ref{tab:benchmark-runtime} \\
C2 & Intended-construct checks: allowance for ER, transfer for AWP & Tables~\ref{tab:alignment-transfer}--\ref{tab:alignment-allowance} \\
C3 & Construct-specific winners (EndorseRank vs.\ AWP) & Table~\ref{tab:alignment-family-ci} \\
C4 & Efficiency--alignment trade-off for EndorseRank, bounded by the degree baseline & Tables~\ref{tab:benchmark-runtime}, \ref{tab:alignment-family-ci}, \ref{tab:holdout-spenders} and~\ref{tab:holdout-diff} \\
C5 & Coupling the two edge types does not beat either single-layer method (bounded negative result) & Tables~\ref{tab:alignment-hybrid}, \ref{tab:tau-diff-hybrid} and~\ref{tab:holdout-diff} \\
\bottomrule
\end{tabular}
}
\end{table}

\dissertationSubheading[sec:research-objectives]{Research objectives}

There are five research objectives, each mapped to RQ1--RQ5. The objective includes the name of the measurement that was used to operationalise the objective, and the criterion that Chapter~\ref{ch:discussion} uses to determine if the objective was met.

\begin{enumerate}
\item[O1.] Find out if EndorseRank and AWP produce different orderings of wallets within the same Arbitrum One cohort (RQ1). Measurement: Kendall $\tau$ between the EndorseRank and AWP scores of the cohort ($n=5{,}521$), with a 95\% confidence interval from bootstrapping. Criterion: the confidence interval does not include $0$ or $1$, where $0$ indicates uncorrelated orderings, and $1$ indicates perfect agreement.
\item[O2.] Assess how well EndorseRank and AWP measure their respective constructs on the same cohort (RQ2). Measurement: family-level Kendall $\tau$ and 95\% confidence intervals for each construct and method, along with paired contrasts $\Delta\tau$ between the two methods on the same construct. Criterion: each method's score on its own construct must have a positive $\Delta\tau$ and confidence interval that does not contain $0$.
\item[O3.] Compare the computation times of PageRank for EndorseRank and AWP when applied to the same number of vertices, and test the effects of node count on the larger set of vertices (RQ3). Measurement: wall-clock solve times (mean and standard deviation of five timed runs, following one warm-up run), memory usage, iteration count, and $|V|$ and $|E|$ for each of the four configurations using a single protocol. Criterion: only report the ratio of the solve times alongside the ratio of $|E|$ so any improvement can be explained by the sparser graph structure.
\item[O4.] Test if scores obtained at $t_1$ can predict approvals at $t_2$ among spenders in the cohort (RQ4). Measurement: Kendall $\tau$ between the scores obtained on 28~February~2026 and the new approval counts between March and May~2026, for the set of $n=1{,}335$ spenders who received funds within that window, with a 95\% confidence interval. Compare the results with those for AWP, and with a paired contrast between EndorseRank and the $t_1$ in-degree. Criterion: EndorseRank's confidence interval does not contain $0$, and the paired difference between EndorseRank and AWP is positive; report the paired difference between EndorseRank and the $t_1$ in-degree with its confidence interval, and judge it on that interval.
\item[O5.] Test if a single PageRank step across both edge types outperforms either single-layer method (RQ5). Measurement: C-PR at $\lambda=0.5$, with $\lambda\in\{0.25,0.75\}$ as a sensitivity test, and S-PR, both evaluated on the same cohort and compared with EndorseRank and AWP through paired confidence intervals on $\Delta\tau$ for the holdout and the same-window labels. Criterion, set prior to the experiment: on the holdout C-PR is as good as EndorseRank on new approval edges and better than AWP on new transfer edges; on the same window it is within the 95\% confidence interval of the best single-layer method for the transfer and allowance labels, and better than both on Sybil stability.
\end{enumerate}

\dissertationSubheading[sec:research-questions]{Research questions}

This research aims to answer three key questions. H1: EndorseRank and AWP result in different wallet orderings for the same wallet set. H2: EndorseRank and AWP are evaluated differently when compared against specific construct checks (allowance check vs transfer check). H3: EndorseRank is faster to compute than AWP at the same value of $n$. We will address these hypotheses through the research questions below that are summarised in Table~\ref{tab:hrq-claim-map}.

\begin{enumerate}
\item[RQ1.] Are the wallet orderings produced by EndorseRank and AWP different, on the same Arbitrum wallet set?
\item[RQ2.] Is EndorseRank more correlated to current allowance checks and current transfer checks than AWP (using Spearman's $\rho$ and Kendall's $\tau$)?
\item[RQ3.] Is EndorseRank faster to compute than AWP on the same wallet set (in terms of time, memory, iterations)?
\item[RQ4.] Once the scores have been computed at time $t_1$, can EndorseRank better predict $t_2$ new owner--spender approvals to inbound spenders than AWP, and better than the allowance degree within the same time window?
\item[RQ5.] Does a coupled PageRank, which operates on both the allowance and the transfer layer, rank wallets better than any single-layer approach (both out-of-sample and in-sample)?
\end{enumerate}

RQ1 tests if there are differences in the centrality measures induced by the allowance and the transfer graph, measured by Kendall's $\tau$ between score vectors for different methods. RQ2 evaluates how well our intended constructs (the transfer local statistics, the allowance local statistics) correlate with different approaches, but since the own-construct proxies are all computed using the same events used to build the graph, we should consider RQ2 as an observation about construct validity, rather than a measure of external validation. RQ3 measures the PageRank computation time, memory usage, and iterations for the same wallet sets, as well as the node-count sensitivity when increasing the wallet pool size. RQ4 is a time-split test of the same construct (i.e.\ endorsements): specifically, predicting future new approvals, not profits from trading. RQ4 replaces the trading-success comparison outlined in the original research proposal which had been withdrawn due to circularity of the outcome-native reference (used as a proxy) with our proposed proxies (see Section~\ref{sub:gf-pr}). RQ5 is a move from testing two alternative edge-swaps to a question of methodology: the coupled operator defined in Section~\ref{sub:coupled-graph} contains both of the single-layer methods as boundary conditions, and the rule defining success for an intermediate point was decided before any scores were calculated (Section~\ref{sec:holdout-protocol}).

\dissertationSubheading[sec:literature-summary]{Summary of the literature review}

Chapter~\ref{ch:literature} contains a review of four main areas of literature; below I summarise what each area has contributed to our approach and where it falls short. First, graph ranking. PageRank \cend{15} calculates node rankings based on recursively propagated endorsements, and is computed via power iteration with per-iteration cost being linear in the number of edges \cend{23}. EigenTrust \cend{16} and TrustRank \cend{17} adapt this approach to the problem of trust in peer-to-peer systems and web spam, respectively, by introducing normalization and seed sets that prevent inflation. These methods provide the solver and complexity arguments employed in this paper, but their edges are limited to links, ratings, or seeds. Second, on-chain reputation. Do and Do\cend{4} apply PageRank and HodgeRank to Ethereum transaction graphs as a social credit metric, while Do, Do, and Nguyen\cend{3} introduce AWP, weighting transfer edges by value and applying a logistic decay over age and validating the ranking against the inbound transfer count and value. Lin et al.\cend{12} propagate account risk through the transaction network. Wu et al.\cend{24} show that graph embeddings of transaction data outperform per-account features for phishing detection. Surveys of blockchain reputation systems \cend{25,26} report that explicit ratings and transfer volume are the most common trust indicators. None of these papers uses the ERC-20 allowance, the on-chain action that logs a deliberate delegation of spending rights \cend{18}, as the graph edge, and the AWP validation compares scores to proxies calculated from the same transfer information used to create the graph. Third, DeFi risk and measurement. Surveys of DeFi \cend{8} and DeFi attacks \cend{27} describe behavioral counterparty risk as one level alongside smart contract, oracle, market, and governance risk. Wash trading \cend{28,29}, extractable value \cend{30}, scam tokens \cend{31}, and airdrop farming \cend{32} all demonstrate that a large amount of transfer activity is artificial, providing the motivation to search for a signal other than payment flow. Fourth, measurement theory and inference. Construct validity \cend{22,33} describes the distinction between what a score measures versus external criteria, while bootstrap resampling \cend{34,35} allows distribution free intervals for rank statistics. The gap that arises (Section~\ref{sec:research-gap}) is a same-cohort comparison of an allowance-edge ranking versus the transfer-edge baseline, with confidence intervals, paired difference tests, a fixed timing protocol reporting edge counts, and a temporal holdout separating same window agreement versus prediction.

\dissertationSubheading[sec:methodology-summary]{Summary of the research methodology}

Chapter~\ref{ch:methodology} employs a quantitative, comparative study design on public ledger data. The Arbitrum One ERC-20 \texttt{Approval} and \texttt{Transfer} event logs from 1~December~2025--31~May~2026 are collected from Google BigQuery; the matched wallet cohort ($n=5{,}521$) includes those addresses with non-zero connectivity in the most recent allowance graph $G_{\text{approve}}$ and the time-decayed transfer graph $G_{\text{transfer}}$. Both graphs are scored using a single power-iteration PageRank implementation ($d=0.85$, $\varepsilon=10^{-8}$, max 300 iterations); EndorseRank uses the sum of latest allowances as edge weights and AWP uses the value weighted, logistically decayed transfers as defined in the baseline paper\cend{3}. The coupled operator C-PR is scored on the same solver with the two edge lists interleaved for each node ($\lambda=0.5$ primary). Rank concordance is measured via Spearman's $\rho$ and Kendall's $\tau$ against three contemporary proxy groups (transfer, allowance, Sybil-adjusted stability), and for RQ4 and RQ5, against new owner--spender approvals and new transfer senders in a three month holdout period following score freezing, alongside the raw $t_1$ degree baselines. All coefficients include a 95\% percentile bootstrap interval generated by 400 paired wallet resamples (seed 42) and pre-specified method comparisons are presented as paired mean intervals. Runtime, peak memory, number of iterations, and the size of $V$ and $E$ are recorded under the same measurement regime (one unrecorded initialization pass followed by five recorded passes, with memoization across runs to ensure identical graphs return identical numbers wherever tables appear). Other sensitivity experiments alter damping, limit to the top tokens, and use an alternative definition of what constitutes a wallet. All pipeline code, configuration files, and machine-readable summaries from which every table was drawn are version-controlled in a public GitHub repository (Section~\ref{sec:reproducibility}), and the experiment uses exclusively public data and does not involve human subjects (Section~\ref{sec:ethics}).

\dissertationSubheading[sec:contributions]{Contributions}

This paper contributes four things to the study of on-chain reputation and DeFi risk. EndorseRank and AWP both apply classic PageRank \cend{15} verbatim to just one type of edge; the coupled operator is the single point in the pipeline where the random walk itself changes. The innovation is thus a new type of edge, a same-cohort contemporaneous comparison, a temporal holdout experiment using degree baselines, and a pre-registered test of whether the two edge types interact.
\begin{itemize}
\item Latest-allowance edge construction (EndorseRank): an endorsement graph in which an edge $u \to v$ represents the current ERC-20 \texttt{approve}/\texttt{permit} allowance held by the spender $v$ from the owner $u$. Since newer approvals override older ones, the latest-allowance edge requires no temporal decay, and it captures an explicit delegation of spending rights rather than a history of payments. No examined reputation method employs this edge (Section~\ref{sec:research-gap}).
\item Same-cohort contemporaneous comparison with statistical inference: EndorseRank and AWP are both computed on a single matched set of wallets, and their outputs are compared on identical transfer, allowance, and Sybil-stability proxies under the AWP validation protocol \cend{3}, using 95\% bootstrap confidence intervals for all coefficients and paired approximate $t$-tests for contrast estimators \cend{34}. Agreement between the two methods is partially tautological. Runtime is included alongside raw edge counts because Claim~C1 follows immediately from sparsity.
\item Temporal holdout experiment with degree baselines: the scores obtained at time $t_1$ are tested on new owner--spender approvals and new transfer senders occurring at time $t_2$ for the spender cohort, alongside the raw number of $t_1$-degree edges that the PageRank smoothing procedure averages over, which separates in-sample agreement from out-of-sample prediction and the propagated PageRank score from its local degree. This, together with the fully open and versioned evaluation pipeline (Section~\ref{sec:reproducibility}; Appendix~\ref{ch:appendix-eval}), supplies the evidence supporting Claims~C1--C5.
\item Coupled operator and its pre-registered evaluation: C-PR blends the allowance and transfer layers for each node using a fixed mixing weight, such that EndorseRank and AWP represent its lower and upper bounds, respectively, and S-PR serves as a seeded ablation. The performance metric, the choice of $\lambda$ and the holdout labels were all set before the experiments were run, and the result is reported unaltered (Claim~C5). For this dataset and cohort, the intermediate values do not exceed the extremes, implying a bound on what any convex combination of the two edge types could achieve, and showing that the advantage of the allowance edge lies in its propagation ability rather than in its ability to seed rankings.
\end{itemize}

Primary numerical results, including cross-method timing, within-family correlations, and runtime, are shown in Chapter~\ref{ch:results} (Tables~\ref{tab:benchmark-runtime} and~\ref{tab:alignment-family-ci}). These results illustrate an alignment-efficiency trade-off: neither reputation method achieves the best score for both edge types.

\dissertationSubheading[sec:scope-limitations]{Scope and limitations}

I study Arbitrum One between 2025-12-01 and 2026-05-31. The main sample consists of $n=5{,}521$ wallet addresses with non-zero endorsement- and transfer-graph connectivity; the number of edges in each graph is reported in Table~\ref{tab:benchmark-runtime}. I report the sensitivity of runtimes to the number of nodes on a larger set of wallets ($N=99{,}080$). I also report the results of my sensitivity analyses for the damping, top-token subgraphs, and the wallet-sample. I do not investigate default labels in unsecured lending. I define constructs and ethics in Chapters~\ref{ch:introduction} and~\ref{ch:methodology}. I make the holdout claim using the spender cohort ($n=1{,}335$). I do not include extended baselines, which are archived for future comparison (Appendix~\ref{ch:appendix-eval}). I do not claim generalizability to other networks than Arbitrum One or outside of my six-month study window. Token-decimal normalization, price oracles, and permit-versus-approve coverage can influence the edge weights. Adversarial manipulation of allowances or transfers could be a potential threat, as I discuss in Chapter~\ref{ch:discussion}. While my Sybil-adjusted proxies help to mitigate the risk of sybil attacks, they do not eliminate it completely \cend{11}.

\dissertationSubheading[sec:roadmap]{Research roadmap}

In Chapter~\ref{ch:literature}, I review graph-based on-chain reputation metrics and ERC-20 allowance semantics. I develop the theoretical framework (hyperlink propagation, domain alignment, and computational complexity) that informs the primary EndorseRank vs.\ AWP comparison.

In Chapter~\ref{ch:methodology}, I describe my data sources, sampling methodology, graph construction for EndorseRank and AWP, contemporaneous construct checks, temporal holdout, computational benchmarks, robustness checks, reproducibility standards, and ethics. I define my matched wallet cohort, spender holdout cohort, the expanded wallet pool I use to scale my system, and the operational procedures to create my Chapter~\ref{ch:results} tables.

In Chapter~\ref{ch:results}, I report my empirical results including my cohort characteristics, runtime and scaling benchmarks (Claim~C1), robustness checks, allowance and transfer alignment (Claims~C2--C3), rank divergence (H1/RQ1), the same-window placement of my coupled operator and the temporal holdout (RQ4, RQ5). I state my primary claims C1--C5 explicitly.

In Chapter~\ref{ch:discussion}, I interpret my findings by research question, develop theoretical and practical implications, formulate a cost model for Sybil attacks on allowance graphs, discuss regulatory and ethical considerations, and evaluate threats to validity and limitations.

In Chapter~\ref{ch:conclusion}, I summarize my contribution to knowledge, answer RQ1--RQ5 sequentially, and suggest avenues for future work, including cross-protocol validation, adversarial robustness, learned coupling weights, and live oracle integration.

Chapter~\ref{ch:literature} begins at the blockchain and DeFi primitives that my allowance construct leverages and concludes at the research gap that this study addresses.
